\section{\code{@defer}, \code{@stream}, and What the Schema Promises}
\label{sec:ch14-defer-and-stream}
\index{defer@\code{@defer}!in a federated graph|(}%

A subscription is not the only way a GraphQL response can arrive in pieces.
\code{@defer} lets a client mark part of a query as lower priority and receive
it in a later chunk; \code{@stream} does the same for the elements of a list.
Both have been in the specification's incremental delivery proposal for years,
and the question a chapter on real time has to answer is whether this stack
does either.

Start with what HotChocolate ships. Both directives exist at 16.6.0, behind two
options:

\begin{minted}[fontsize=\footnotesize,breaklines]{csharp}
public bool EnableDefer { get; set; }
public bool EnableStream { get; set; }
\end{minted}

\index{EnableDefer@\code{EnableDefer}!off by default}%
Auto-properties with no initialiser, so both are false. That is worth reading
against \code{public bool EnableTag \{ get; set; \} = true;}, six lines further
down the same file, which is what an option that is meant to be on looks like
there. The interface documents both as preview features. Mosaic leaves
them off, and introspecting any subgraph confirms it: the directive list is
\code{skip}, \code{include}, \code{deprecated}, \code{shareable}, \code{key},
\code{specifiedBy}, \code{link}, and sending \code{@defer} to a subgraph
directly gets a flat refusal.

\begin{minted}[fontsize=\footnotesize,breaklines]{text}
The specified directive `defer` is not supported by the current schema.
\end{minted}

\subsection*{The router says otherwise}

\index{defer@\code{@defer}!added by the router}%
Now introspect the router. Its client-facing schema advertises six directives,
and two of them appear in no subgraph:

\begin{minted}[fontsize=\footnotesize,breaklines]{text}
include  skip  deprecated  specifiedBy  oneOf  defer
\end{minted}

\code{oneOf} is uninteresting: it describes input objects and the router adds
it because the specification has it. \code{defer} is the one to look at. Its
locations are \code{FRAGMENT\_SPREAD} and \code{INLINE\_FRAGMENT}, and it is in
the client schema of a graph composed entirely from subgraphs that have never
heard of it. The router puts it there. So a client generator pointed at this
graph, or a developer reading the schema in a playground, is told the feature
is available.

Use it, on a fragment whose fields live in a different subgraph from the rest
of the query, and ask for the multipart response an incremental delivery would
need:

\begin{minted}[fontsize=\footnotesize,breaklines]{text}
POST /graphql   Accept: multipart/mixed

{ browseProducts(first: 1) { nodes { title ... @defer { price { amount } } } } }
\end{minted}

\begin{minted}[fontsize=\footnotesize,breaklines]{text}
HTTP/1.1 200 OK
Content-Type: application/json; charset=utf-8

{"data":{"browseProducts":{"nodes":[{"title":"Beech Cutting Board",
  "price":{"amount":69.00}}]}}}
\end{minted}

\index{defer@\code{@defer}!accepted and ignored}%
One response, \code{application/json} rather than multipart, and the deferred
field sitting in it as though the directive were not there. The directive is
accepted by validation and then has no effect. \code{@stream} at least says so:

\begin{minted}[fontsize=\footnotesize,breaklines]{text}
{"errors":[{"message":"directive: stream undefined","path":["query","products"]}]}
\end{minted}

Which of those two is the better failure is not a close call. \code{@stream} is
rejected in a way a client sees during development, once. \code{@defer} is
advertised, accepted, silently flattened, and will look like it works right up
until somebody measures a page that was supposed to paint in two stages and
finds it painting in one.

\subsection*{Why nobody is in a hurry to fix this}

\index{incremental delivery!vendor positions}%
It would be easy to file this as a gap and move on, and I do not think that is
what it is. Jens Neuse, WunderGraph's chief executive and co-founder, published
a post in 2023 arguing that \enquote{the @defer and @stream directives are
complete overkill and the whole idea of incrementally loading data can be done
in a much simpler way}~\autocite{neuse2023defer}. The company that writes this
router does not think the feature is worth its complexity, and has said so
under a name. HotChocolate's position is quieter and points the same way: both
options exist, both are off, and its own documentation comment calls them
preview features.

You can disagree with that. What you cannot do is read the router's schema and
conclude anything about it, which is the part I would push back on. A directive
in a published schema is a promise, and this is a graph making one it does not
keep. If the answer is that incremental delivery is not worth doing, the
consistent version of that answer is not to advertise the directive.

\index{incremental delivery!what to do instead}%
The practical advice is short. If you want part of a page later, ask for it
later: a second query is a thing every layer in this stack understands, every
gate in this book can assert, and no client library has to negotiate a
multipart parser for. If you want data pushed as it changes, the previous five
sections are the mechanism this stack actually implements, and one of them does
not even need the service to be running.

\medskip
Chapter~\ref{ch:hard-modeling-problems} ended by saying a federated graph can
make the shapes agree and cannot make the meanings agree. This chapter adds a
smaller one to that: it cannot make a schema honest either, and both of the
things it got wrong here were invisible from every schema, every log and every
composition check. Which is a good place to stop, because the next chapter is
about who is allowed to ask, and a connection that stays open for an hour is
the hardest case that question has.

\index{defer@\code{@defer}!in a federated graph|)}%
